\section{Evaluation Methodology}
\label{sec:setup}

The evaluation is organized around four answerable questions. RQ1 asks whether
an ASK controller selects the correct write action under controlled evidence.
RQ2 asks whether a transition preserves state and lineage invariants. RQ3 asks
how false commits, missed updates, and unnecessary abstentions change under
different costs. RQ4 asks whether governed writes improve downstream answers
and whether later failures can be attributed to writing, retrieval, or
generation. Generalization to independently collected natural dialogue is not
an answered research question in this version.

\begin{center}
\small
\begin{tabular}{@{}p{0.22\linewidth}p{0.22\linewidth}p{0.46\linewidth}@{}}
\toprule
Evidence stream & Status & Supported use \\
\midrule
Reference protocol tests & Primary conformance & Transition and invariant semantics, not semantic action accuracy \\
\mgpc{} & Primary controlled & RQ1--RQ3 on a controlled synthetic single-slot diagnostic \\
\lme{} V12 & Downstream final candidate & RQ4 on one fixed slice; no per-candidate action gold \\
\mgp{} & Historical diagnostic & Construct-failure discovery and mechanism examples only \\
External adapters & Mapped diagnostic & Harness compatibility; not native-system capability \\
MGP-Natural v0.1.0/v0.2.2 & Quarantined, not scored & Feasibility and construction evidence only; no independent gold or headline metric \\
\bottomrule
\end{tabular}
\end{center}

\begin{figure}[!h]
\centering
\fbox{\begin{minipage}{0.94\linewidth}
\centering
candidate and evidence
$\rightarrow$ \textbf{action} (accuracy, class recall, false commit)\\[2pt]
$\rightarrow$ \textbf{state and lineage} (invalid state, retained predecessor,
active successor)\\[2pt]
$\rightarrow$ \textbf{retrieval} (active-only recall, stale retrieval)\\[2pt]
$\rightarrow$ \textbf{answer} (task accuracy, refusal, generation error)
\end{minipage}}
\caption{Layered evaluation in ASK. A later correct answer does not erase an
incorrect write, and a correct write is not charged for a later retrieval or
generation failure.}
\label{fig:evaluation-layers}
\end{figure}

\subsection{Historical v2 systems}

The main test comparison contains the transparent reference policy, the
structured LLM policy, and the unstructured zero-shot policy described in
Section~\ref{sec:method}. Before interpreting them, we run type-only,
conflict-type-only, and source/reliability/temporal lookup probes fitted on
train+dev. Four structurally limited systems (no memory, flat retrieval,
recency, and naive overwrite) and component ablations are retained as pilot
mechanism checks in the appendix; they are not competitive peers.

A six-demonstration LLM decision run produced $174/180$ predictions. Aligning
its six missing identifiers to \texttt{INVALID} yields a coverage-aware audit,
but the run is incomplete and excluded from every comparison and hypothesis
test.

Adapter runs for Mem0, Letta, Graphiti, and MemOS are preserved as harness
diagnostics in the artifact. Their native operations do not match our action
vocabulary, so mapped action accuracies are excluded from the main empirical
claim. A successor evaluation must compare observable final state, stale
retrieval, downstream answer, latency, and cost under each system's native API.

\subsection{Historical v2 metrics and denominators}

Primary decision-level metrics are computed over the held-out $90$-item test:
action accuracy, macro-F1, and per-action precision/recall/F1. We report
action false-commit rate over gold-\abstain{} items, missed-update rate over
gold-\supersede{} items, and unnecessary-abstention rate over non-\abstain{}
items. A zero event count is accompanied by a finite-sample upper bound.

Memory-state outcomes are lineage integrity (a five-condition per-item score on
gold-\supersede{} cases) and stale-retrieval rate. We additionally report
top-label ECE/Brier, risk--coverage, and expected loss under balanced,
commit-averse, and update-responsive cost matrices. Confidence results are
descriptive because neither policy was calibrated on an independent set. Full
formulas and commands are in Appendix~\ref{app:metrics}.

\subsection{Historical v2 statistical methods}
\label{sec:stats}

Systems are paired by scenario id. The primary uncertainty estimate is a
$10{,}000$-replicate bootstrap over \texttt{template\_family\_id}; we repeat it
with minimal-pair groups plus singleton unpaired items. We report family-macro
and item-weighted estimates. Exact two-sided McNemar tests are secondary because
item independence is doubtful. No equivalence or non-inferiority claim is made:
no margin was pre-specified.

For zero observed action false commits we report the exact one-sided 95\% upper bound
$1-0.05^{1/n}$ rather than interpreting the observed rate as population zero.
These calculations quantify uncertainty inside this collection; they do not
repair its construct leakage or establish external generalization.

Annotation agreement is reported as Cohen's $\kappa$ for one independent human
annotator against the author gold and, separately, Fleiss $\kappa$ for a
four-model audit panel. The latter is never described as human
inter-annotator agreement.

\subsection{\mgpc{} systems, metrics, and statistics}

The \mgpc{} comparison covers sixteen configurations on the public development
split: the frozen controller, the successor controller in both
its as-run and guard-composed forms, four structural baselines (no memory, flat
retrieval, recency, naive overwrite), an unstructured zero-shot decision policy,
the deterministic five-feature scorecard through its v4.3 adapter, and four
component ablations (no abstention, no provenance, no typing, no lineage).
Three external memory systems, Letta, Mem0, and Graphiti, were run through adapters
we wrote. Because their native operations do not match our action vocabulary,
their scores depend on an action mapping and are reported as harness
diagnostics, not as capability measurements of those systems in their intended
use. All model-backed rows use \texttt{qwen/qwen3.5-9b} at temperature $0$,
seed $4301$, $128$ maximum new tokens, and concurrency $1$.

Primary decision metrics are action accuracy, macro-F1, per-action
precision/recall/F1, Cohen's $\kappa$, and the multiclass Matthews correlation;
the last two are informative here precisely because the splits are exactly
balanced. We decompose errors by operational harm rather than reporting a single
error rate: action false commit on gold-\abstain{} items, separated into the
state-preserving \abstain$\to$\keep{} and state-corrupting
\abstain$\to$\supersede{} components; missed update on gold-\supersede{} items;
destructive overwrite of a gold-\keep{} record; and unnecessary abstention over
items that required a decision. Every rate carries a Wilson $95\%$ interval, and
a zero count is reported with a finite-sample upper bound rather than as zero
risk. Expected loss is computed under balanced, commit-averse, and
update-responsive cost matrices fixed before scoring.

Hidden-split uncertainty uses a $5{,}000$-replicate component-cluster bootstrap
at seed $4301$, because items within a semantic component share a scenario
skeleton. On the public development split, component-cluster resampling is
unavailable because the component map is private and the public scorer does not
read it. Those intervals are therefore $10{,}000$-replicate item bootstraps and are
narrower than a component-aware estimate would be. Paired system contrasts on
development use exact two-sided McNemar tests with Holm correction over the
reported family; they inherit the same independence assumption and are
interpreted accordingly (Section~\ref{sec:limitations}).

Operational accounting is taken from the per-response records rather than
from wall-clock timing of the harness: provider-reported latency per response
(mean, median, p95, maximum), prompt and completion tokens, transport failures,
in-client retries, and invalid-response counts. Monetary cost is reported only
where the provider emitted a price field; on every run in this paper it did
not, and we record cost as unavailable rather than substituting list-price
arithmetic.

\subsection{\mgpc{} staged execution and reproducibility}

The staged protocol and its one-shot hidden discipline are specified in
Section~\ref{sec:v43contract}. In summary: train supports engineering checks,
development carries every prompt, configuration, guard, ablation, cost, and
latency choice, and the sealed hidden split was opened exactly once, after the
configuration was frozen and hashed and the prediction file was
coverage-validated and sealed. The direct methodological consequence, which
governs how Section~\ref{sec:results} may be read, is that no hidden baseline,
ablation, or paired comparison exists.
